\section{Viable Prefix dan Konflik}
\label{sec:ch11-viable}

\subsection{Viable Prefix}

\begin{definisi}[Viable prefix]
Viable prefix adalah prefiks dari bentuk kalimat yang dapat muncul
pada stack parser shift-reduce tanpa melewati handle.
\end{definisi}

Viable prefix menentukan apakah stack masih dapat diperluas menjadi
bentuk kalimat yang valid. Jika stack bukan viable prefix, parser
tidak mungkin mencapai simbol awal.

\subsection{Konflik Shift-Reduce}

Konflik shift-reduce terjadi ketika parser dapat memilih antara
shift dan reduce untuk stack dan token yang sama. Konflik ini sering
muncul pada grammar ambigu, misalnya dangling else atau ekspresi
tanpa tingkat prioritas.

\begin{contoh}[Konflik shift-reduce]
Pada grammar
\[
E \to E + E \mid E * E \mid \code{id},
\]
setelah stack berisi \(E + E\) dan token berikutnya \(*\), parser
dapat mengurangi \(E + E\) menjadi \(E\) atau menggeser \(*\)
untuk membentuk \(E * E\). Pilihan ini menentukan apakah
penjumlahan atau perkalian yang lebih dulu dikelompokkan.
\end{contoh}

\subsection{Konflik Reduce-Reduce}

Konflik reduce-reduce terjadi ketika dua aturan berbeda dapat
diterapkan pada handle yang sama. Konflik ini biasanya menunjukkan
bahwa grammar terlalu ambigu atau memiliki aturan yang tumpang
tindih.

\begin{figure}[H]
\centering
\begin{tikzpicture}[node distance=6mm and 12mm]
  \node[decision, text width=2.8cm] (choice) {Stack dan\\token};
  \node[phase, text width=2.8cm, below left=of choice] (shift)
        {Shift};
  \node[phase, text width=2.8cm, below=of choice] (reduce)
        {Reduce};
  \node[phase, text width=2.8cm, below right=of choice] (error)
        {Error};
  \draw[arrow] (choice) -- (shift);
  \draw[arrow] (choice) -- (reduce);
  \draw[arrow] (choice) -- (error);
\end{tikzpicture}
\caption{Tiga kemungkinan aksi pada parser shift-reduce}
\label{fig:ch11-actions}
\end{figure}

\subsection{Menyelesaikan Konflik}

Konflik dapat diselesaikan dengan:
\begin{itemize}[leftmargin=*]
  \item mengubah grammar (precedence, associativity, factoring);
  \item menambahkan aturan prioritas pada parser generator;
  \item memilih shift atau reduce berdasarkan kebijakan tetap.
\end{itemize}

Pemilihan kebijakan harus didokumentasikan. Parser generator seperti
Bison secara default memilih shift, tetapi keputusan ini dapat
menyembunyikan bug jika tidak disengaja.
